\subsection{QEMU, The Quick Emulator}
\label{chap:Back_QEMU}

The core of this project lies in the hardware emulation platform. QEMU was 
selected over other options after meeting several requirements revealed in previous 
analyses, which will be explained later in this document. This section of the chapter 
aims to introduce QEMU as a tool, A in-depth analysis of the tool will be done in 
the next chapter XXXX.

The Quick Emulator or as it is mostly known QEMU, is a generic and 
open source machine emulator and virtualizer that allows the user operating 
systems and programs designed for one machine architecture on a completely different 
architecture. Created by Fabrice Bellard, QEMU has evolved into a sophisticated 
virtualization solution that supports a wide range of CPU architectures including 
x86, ARM, RISC-V, PowerPC, and MIPS. 

QEMU operates through Dynamic Binary Translation (DBT), based on JIT 
compilation, to achieve fast emulation speed. The use of DBT enables running 
binary code generated for a target (or a guest) processor on another host processor, 
which has an entirely different ISA. This means that it translates the binary code 
associated to the target CPU into the instructions of the host CPU. Then, it executes 
the translated instructions on the host CPU.

Since its release in 2005, many features have been added to QEMU, such as 
peripheral support, support for different architectures, virtualisation and performance 
improvements. Nowadays, QEMU offers multiple modes of operation:

\begin{enumerate}
  \item \textbf{Full-system emulation:} 
  Complete virtual machine emulation with virtual 
  peripherals. The QEMU's default translation engine is DBT or parallelized-DBT, 
  called Multi-Threaded Tiny Code Generator (MTTCG)

  \item \textbf{User-mode emulation:} 
  Running individual applications compiled for different 
  architectures than the guest system. In this mode QEMU emulates only 
  the CPUs. It executes guest instructions, captures system calls, and sends 
  them to the host kernel. 
  
  \item \textbf{Hardware-virtualization:} 
  QEMU use hardware extensions like KVM or Xen 
  to emulate the guest code, with near-native performance. In this case, QEMU 
  deals with all virtualisation aspect that can not be emulated (peripherals or 
  input/output management). The guest code execution is a task of the host 
  virtualisation support. The virtualisation requires host hardware support to 
  execute the guest code directly, so it only works if the guest has the same 
  instruction set architecture as the host.
\end{enumerate}


%==== Full-System ====%
\subsubsection{Full-System Emulation}
\label{chap:Back_QEMU_FSE}

For the current project, and because the necessity to emulates an entire development 
board, including the STM32F405 microcontroller, the memory, the peripherals, 
the buses and firmware, QEMU will be used in its full system emulation mode.
Regardless of the mode, at the highest level QEMU operates as a conventional userspace 
application on the host operating system,\ref{fig:full_system_emulation}. Invoked in full-system mode, the 
main QEMU process performs an initialization of all the virtual componentes of the 
hardware, and then enters a main loop in which it repeatedly executes blocks of 
guest code and handles peripheral events.

\begin{figure}[htbp]
    \centering
%    \includegraphics[width=0.7\textwidth]{Resources/Draw/Full-system.pdf}
    \caption{Full system emulation architecture.}
    \label{fig:full_system_emulation}
\end{figure}